\documentclass[11pt]{article}

\usepackage[margin=1in]{geometry}
\usepackage{amsmath, amssymb}
\usepackage{mathtools}
\usepackage{booktabs}
\usepackage{array}
\usepackage{float}
\usepackage[colorlinks=true, linkcolor=blue, citecolor=blue, urlcolor=blue]{hyperref}

\newcommand{\ii}{\mathrm{i}}
\newcommand{\dd}{\mathrm{d}}
\newcommand{\PhiS}{\Phi^{\mathrm{S}}}
\newcommand{\Kmn}[2]{K_{#1}^{(#2)}}
\newcommand{\codeus}{\textunderscore\allowbreak}
\newcommand{\codehy}{-\allowbreak}
\newcommand{\code}[1]{{\ttfamily\let\_\codeus\let\-\codehy #1}}
\emergencystretch=2em

\title{Analytic azimuthal modes of the equatorial effective source}
\author{EffectiveSource project notes}
\date{2026-10-17}

\begin{document}
\maketitle

\begin{abstract}
We describe a method to compute the azimuthal Fourier modes of the effective
source of a scalar charge on an equatorial Kerr geodesic, together with the
modes of the singular field and of its first spatial derivatives, without any
quadrature in the azimuth. At a fixed field point and particle position, each of
these quantities can be written as a single fraction: an even polynomial in
$u=\sin^2(\psi/2)$ plus $\sin\psi$ times another polynomial, divided by a fixed
half-integer power of the quadratic distance $s^2=\alpha+\beta u$, where $\psi$
is the shifted azimuthal separation from the particle. The azimuthal integral of
such a fraction is a finite linear combination of the kernels
$\int_0^{2\pi}\cos(\mu\psi)\,(\alpha+\beta\sin^2\tfrac{\psi}{2})^{-\nu}\,\dd\psi$,
which follow from two complete elliptic integrals and two three-term
recurrences. Near the particle the numerator coefficients are the small
remainder of large cancelling terms, and we compute them, the kernels and the
contractions in double-double arithmetic where this cancellation requires it.
Part~\ref{part:theory} develops the method and Part~\ref{part:implementation}
describes its implementation, its tests and its run time. The modes agree with
high-precision quadrature of the same functions to better than $10^{-13}$
relative to each mode.
\end{abstract}

\tableofcontents

\part{Theory}
\label{part:theory}

\section{Introduction}
\label{sec:intro}

In the effective-source approach to the self-force, the retarded field of a
point charge is split into a singular field $\PhiS$, known in closed form near
the worldline, and a regular remainder. The remainder satisfies the wave
equation with the effective source $S=-\Box\PhiS$, which is bounded at the
particle and can be handled by a standard field solver. For a particle on an
equatorial orbit of a Kerr black hole, the background is stationary and
axisymmetric, and it is natural to decompose every field into azimuthal modes
\begin{equation}
  \label{eq:mode}
  f_m(r,\theta) = \int_0^{2\pi} f(r,\theta,\phi)\, e^{-\ii m\phi}\,\dd\phi ,
\end{equation}
evaluated at fixed time $t$. A frequency-domain solver needs $S_m$ at every
point of its grid and, for eccentric orbits, at many times along the orbit. On
the boundary of the region where the effective source is used, it also needs
the modes of $\PhiS$, $\partial_r\PhiS$ and $\partial_\theta\PhiS$.

The integrand of Eq.~\eqref{eq:mode} is sharply peaked in $\phi$ when the field
point is close to the orbit, and the peak narrows as the field point approaches
the particle. A quadrature rule in $\phi$ must resolve the peak, and the number
of nodes it needs grows without bound as the separation decreases. A second
difficulty is cancellation. The source is bounded at the particle, but it is
the sum of second derivatives of $\PhiS$, each of which grows like the inverse
cube of the distance. Any method that integrates these terms separately and
adds the results has its relative error amplified by the ratio of the terms to
their sum.

This note describes a method that avoids both problems. We first show that, at
a fixed field point, every quantity we need is a single rational function of
the azimuth with a fixed denominator, in which all the cancellation between the
singular terms has already taken place in a small set of numerator
coefficients (Sec.~\ref{sec:rational}). We then show that the azimuthal
integral of such a function is a short linear combination of one family of
kernels (Sec.~\ref{sec:integration}), and that these kernels can be computed
to full precision from elliptic integrals and recurrences
(Sec.~\ref{sec:kernels}). Section~\ref{sec:precision} discusses the precision
that each step requires. We begin with the form of the singular field and the
source in Sec.~\ref{sec:field}.

\section{The singular field and the effective source}
\label{sec:field}

\subsection{Local coordinates and the distance function}

We use Boyer--Lindquist coordinates $(t,r,\theta,\phi)$ and place the particle
at $(r_p,\pi/2,\phi_p)$ at time $t$. The separations of the field point from
the particle are $\delta r=r-r_p$, $\delta\theta=\theta-\pi/2$ and
$\delta\phi=\phi-\phi_p$. The singular field is written in terms of a shifted
azimuthal separation and its half-angle sine,
\begin{equation}
  \label{eq:psi}
  \psi = \delta\phi - c\,\delta r, \qquad \sigma=\sin\frac{\psi}{2}, \qquad u=\sigma^2,
\end{equation}
where the constant $c$ depends on the state of the particle. The singular field
depends on the azimuth only through $u$ and $\sin\psi$. Both are
$2\pi$-periodic in $\phi$, as a field must be, and near the particle they agree
with $\psi^2/4$ and $\psi$.

The building block of the singular field is the quadratic distance function
\begin{equation}
  \label{eq:s2}
  s^2 = \alpha + \beta u, \qquad \alpha = \alpha_{20}\,\delta r^2 + \alpha_{02}\,\delta\theta^2,
\end{equation}
where $\alpha_{20}$, $\alpha_{02}$ and $\beta$ are positive constants that,
like $c$, depend only on the state of the particle. At leading order $s^2$ is
the squared distance from the particle in the local metric. For a given field
point $(r,\theta)$, the quantity $\alpha$ is the smallest value of $s^2$ over
all azimuths, reached at $\psi=0$. We measure the separation of the field point
from the particle by the dimensionless ratio
\begin{equation}
  \label{eq:C1}
  C_1 = \frac{\alpha}{\beta},
\end{equation}
which vanishes only on the orbit, $\delta r=\delta\theta=0$. Near $\psi=0$ we
have $s^2\approx\alpha+\beta\psi^2/4$, so the azimuthal width of the peak of
any power of $s^{-2}$ is about $2\sqrt{C_1}$.

\subsection{The singular field}

The singular field is the ratio of a polynomial to the seventh power of $s$,
\begin{equation}
  \label{eq:puncture}
  \PhiS = \frac{A}{s^7}, \qquad
  A = \sum_{ijk} A_{ijk0}\,\delta r^i\,\delta\theta^j\,\sigma^k
    + \sin\psi\sum_{ijk} A_{ijk1}\,\delta r^i\,\delta\theta^j\,\sigma^k .
\end{equation}
The coefficients $A_{ijkl}$ depend only on the state of the particle. Only even
powers $j$ of $\delta\theta$ appear, because the problem is symmetric under
reflection through the equatorial plane, and only even powers $k$ of $\sigma$.
Odd behaviour in $\psi$ enters through one overall factor of $\sin\psi$. Every term
has total degree $6\le i+j+k+l\le9$, counting $\sin\psi$ as degree one, and
there are 85 coefficients in total. Since $A$ is at least of degree six and
$s^7$ is of degree seven, the singular field behaves like $1/s$ at leading
order, as the field of a point charge must. The higher-degree terms are the
corrections that make the effective source bounded.

For an eccentric orbit the particle moves, and the coefficients $A_{ijkl}$,
$\alpha_{20}$, $\alpha_{02}$, $\beta$ and $c$ change along the orbit. The time
derivatives of the singular field therefore also require their first and second
derivatives along the orbit, $\dot A_{ijkl}$ and $\ddot A_{ijkl}$, and those of
the other constants, together with the velocity and acceleration of the
particle. We refer to all of these numbers as the particle data.

\subsection{Derivatives of the singular field}

The singular field depends on $t$, $r$ and $\phi$ through $\delta r$, through
$\psi$ and through the time dependence of the particle data, and on $\theta$
only through $\delta\theta$. The dependence through $\psi$ follows from
Eq.~\eqref{eq:psi},
\begin{equation}
  \label{eq:dpsi}
  \partial_\phi\psi=1, \qquad \partial_r\psi=-c, \qquad
  \partial_t\psi = c\,\dot r_p - \dot c\,\delta r - \dot\phi_p ,
\end{equation}
and the derivatives of $\sigma$ and $\sin\psi$ follow by the chain rule, for
example $\partial_r\sigma=-\tfrac{c}{2}\cos(\psi/2)$. Differentiating
Eq.~\eqref{eq:puncture} gives, for any coordinates $X$ and $Y$,
\begin{align}
  \partial_X\PhiS &= \frac{2s^2\,\partial_X A - 7A\,\partial_X s^2}{2s^9}, \label{eq:first}\\
  \partial_X\partial_Y\PhiS &= \frac{63A\,\partial_X s^2\,\partial_Y s^2
    - 14s^2\bigl(\partial_X A\,\partial_Y s^2 + \partial_Y A\,\partial_X s^2 + A\,\partial_X\partial_Y s^2\bigr)
    + 4s^4\,\partial_X\partial_Y A}{4s^{11}}, \label{eq:second}
\end{align}
where every derivative is a total derivative that includes the dependence
through $\psi$. Each derivative lowers the power of $s$ by two. The numerators
are polynomials in the separations, in $\sigma$, $\cos(\psi/2)$ and $\sin\psi$,
and in the particle data.

\subsection{The effective source}

The effective source is the scalar wave operator of the Kerr metric applied to
the singular field. In Boyer--Lindquist coordinates it takes the form
\begin{equation}
  \label{eq:source}
  S = -\Box\PhiS = \sum_X b_X(r,\theta)\,\partial_X\PhiS
    + \sum_{X,Y} b_{XY}(r,\theta)\,\partial_X\partial_Y\PhiS,
\end{equation}
where the first sum runs over $\partial_r$ and $\partial_\theta$ and the second
over $\partial_r^2$, $\partial_\theta^2$, $\partial_t^2$, $\partial_t\partial_\phi$
and $\partial_\phi^2$. The operator coefficients $b_X$ and $b_{XY}$ depend on
the field point only, because the background is stationary and axisymmetric.

Equation~\eqref{eq:second} shows that each term of Eq.~\eqref{eq:source}
grows like $s^{-3}$ near the particle. The singular field is constructed so
that these terms cancel to the order of the expansion, and their sum is
bounded. The size of the individual terms relative to the sum is the origin of
the cancellation discussed in the introduction.

\subsection{Quadrature in the azimuth}
\label{sec:direct}

The integrand of Eq.~\eqref{eq:mode} is analytic in $\phi$, with singularities
in the complex plane where $s^2$ vanishes. Setting $\alpha+\beta\sin^2(\psi/2)=0$
gives $\psi=\pm2\ii\,\mathrm{arsinh}\sqrt{C_1}$, so the singularities lie a
distance of about $2\sqrt{C_1}$ from the real axis. The trapezoidal rule with
$N$ nodes then has an error that decreases like
$\exp(-2N\,\mathrm{arsinh}\sqrt{C_1})$, and reaching a relative error of
$10^{-14}$ requires $N\approx16/\sqrt{C_1}$ nodes, about $5\times10^5$ at
$C_1=10^{-9}$. The same distance sets the decay of the modes with $m$, which we
meet again in Sec.~\ref{sec:mu}. Quadrature also leaves the cancellation in
the source untouched, since the integrand at each node is the sum of
Eq.~\eqref{eq:source}.

\subsection{Term-by-term closed forms}
\label{sec:termwise}

An earlier approach avoids quadrature by integrating each term of the singular
field in closed form and by obtaining the derivatives and the source from these
closed forms. We describe it in some detail, because the ways in which it fails
in double precision motivate each step of the method of
Secs.~\ref{sec:rational}--\ref{sec:precision}.

\paragraph{Construction.} The polynomial $A$ of Eq.~\eqref{eq:puncture} is
first written as
\begin{equation}
  \label{eq:A_split}
  A = \sum_{j=0}^{4}\bigl(R_j + \sin\psi\,I_j\bigr)\,\sigma^{2j},
\end{equation}
where the ten coefficients $R_j$ and $I_j$ are polynomials in $\delta r$ and
$\delta\theta$. The mode of the singular field then requires the integrals
\begin{equation}
  \label{eq:old_integrals}
  c_{mj} = \int_0^{2\pi}\frac{\cos(m\psi)\,\sigma^{2j}}{s^7}\,\dd\psi, \qquad
  d_{mj} = \int_0^{2\pi}\frac{\sin(m\psi)\,\sin\psi\;\sigma^{2j}}{s^7}\,\dd\psi,
\end{equation}
in terms of which
$(\PhiS)_m = e^{-\ii m\phi_0}\sum_j\bigl(R_j\,c_{mj} - \ii\,I_j\,d_{mj}\bigr)$,
with the constant $\phi_0=\phi_p+c\,\delta r$. Each integral has a closed form
in the complete elliptic integrals $K(k)$ and $E(k)$ with
$k^2=\beta/(\alpha+\beta)=1/(1+C_1)$,
\begin{equation}
  \label{eq:old_closed}
  c_{mj} = \frac{4\bigl[p_{mj}(C_1)\,K(k) + q_{mj}(C_1)\,E(k)\bigr]}{\beta^{7/2}\,C_1^3\,(1+C_1)^{5/2}},
\end{equation}
where $p_{mj}$ and $q_{mj}$ are polynomials of degree up to $m+2+j$ with
rational coefficients. The integrals $d_{mj}$ have the same form, with a
denominator proportional to $C_1^2(1+C_1)^{3/2}$. The polynomials are derived
once by computer algebra and tabulated for each $m$.

The modes of the derivatives follow from the closed forms. At fixed $\phi$ the
azimuthal integral commutes with $\partial_t$, $\partial_r$ and
$\partial_\theta$, so $(\partial_X\PhiS)_m=\partial_X(\PhiS)_m$ for these
coordinates, and $(\partial_\phi\PhiS)_m=\ii m\,(\PhiS)_m$. The closed form
depends on $(t,r,\theta)$ through $R_j$ and $I_j$, through $C_1$, which enters
$K$, $E$, the polynomials and the denominators, and through the phase
$e^{-\ii m\phi_0}$. The chain rule uses
\begin{equation}
  \label{eq:dKdC}
  \frac{\dd K}{\dd C_1} = \frac{C_1K-(1+C_1)E}{2C_1(1+C_1)}, \qquad
  \frac{\dd E}{\dd C_1} = \frac{K-E}{2(1+C_1)},
\end{equation}
and their derivatives, and the modes of the second derivatives follow by
differentiating once more. The mode of the source is finally obtained by
applying Eq.~\eqref{eq:source} to these modes, with $\partial_\phi$ replaced
by $\ii m$.

For the singular field, this route and the method of
Secs.~\ref{sec:rational}--\ref{sec:kernels} compute the same integrals. The
kernels of Sec.~\ref{sec:integration} are themselves combinations of $K$ and
$E$ with coefficients rational in $C_1$, generated by the recurrences of
Sec.~\ref{sec:kernels}, so Eq.~\eqref{eq:old_closed} and the contractions of
Sec.~\ref{sec:integration} are two ways of writing the same exact result. The
same holds for the first derivatives of the singular field. Because the
azimuthal integral commutes with the derivatives, differentiating the closed
forms after the projection gives the same modes as differentiating the singular
field before it, as in Sec.~\ref{sec:fields}. The two routes differ in how the
exact expressions are evaluated in floating point and in how the source is
assembled, and in double precision the term-by-term route fails in three
ways.

\paragraph{Cancellation at large separation and high $m$.} The closed form of
mode $m$ is, algebraically, what the recurrence in $\mu$ of Sec.~\ref{sec:mu}
produces when it is applied $m$ times to $K$ and $E$. Its polynomials grow
rapidly with $m$, and the largest tabulated coefficient is about $9$ for
$m=0$, $4\times10^2$ for $m=2$, $4\times10^9$ for $m=10$ and $5\times10^{17}$
for $m=20$. The mode itself decreases like $\rho^{-m}$, with $\rho$ given by
Eq.~\eqref{eq:rho}. The terms of the closed form therefore cancel, and the
rounding errors of $K$, $E$ and the coefficients are amplified by roughly
$\rho^{2m}$, the same factor as for the forward recurrence. At $C_1=1$, where
$\rho\approx5.8$, this factor exceeds $10^{16}$ for $m\ge11$, and the high modes
have no correct digits.

\paragraph{The elliptic integrals near $k=1$.} Near the particle $k\to1$. If $K$
and $E$ are computed from $k^2=1/(1+C_1)$, the rounding of $k^2$ gives the small
complement $1-k^2\approx C_1$ a relative error of order $\epsilon/C_1$, where
$\epsilon\approx10^{-16}$ is the unit roundoff. Since
$K(k)\approx\tfrac12\ln\bigl(16/(1-k^2)\bigr)$, the absolute error of $K$ is
about $\epsilon/(2C_1)$, which limits the modes of $\PhiS$ to a relative
accuracy of about $10^{-10}$ at $C_1=10^{-8}$. Computing $K$ and $E$ from the
complementary parameter, as in Sec.~\ref{sec:kernels}, removes this error.

\paragraph{Cancellation in the source.} The modes of the second derivatives
grow like $1/C_1$ near the particle, while the mode of the source remains
finite. Near $\psi=0$ the azimuthal integral of $s^{-3}$ is
$\int(\alpha+\beta\psi^2/4)^{-3/2}\,\dd\psi=4/(\alpha\sqrt\beta)$, which sets
this scaling. When the source is assembled from these modes by
Eq.~\eqref{eq:source}, the separately projected terms cancel to the same
degree, so the relative error of $S_m$ grows like $1/C_1$ as the field point
approaches the particle, however $K$ and $E$ are evaluated. The singular field
and its first derivatives involve no such cancellation, and their modes from the
differentiated closed forms are as accurate as those of $\PhiS$ itself.

In addition, the tables limit the construction to the modes for which closed
forms were derived, and extending them to higher $m$ increases the
cancellation of the first kind. Section~\ref{sec:old} quantifies each of these
errors.

The method of Secs.~\ref{sec:rational}--\ref{sec:precision} addresses each
failure. It combines the terms of the source before the azimuthal integration,
so that the cancellation of the third kind takes place in a small number of
numerator coefficients that can be computed in higher precision. It computes
the kernels by the recurrence in the direction in which that recurrence is
stable, or in a basis in which the high modes do not cancel, which avoids the
first kind. It evaluates $K$ and $E$ from the complementary parameter, which
avoids the second, and its kernels are available to high $m$ without
tabulation.

\section{A single fraction for each quantity}
\label{sec:rational}

\subsection{Dependence on the azimuth at a fixed field point}

The method rests on the observation that, at a fixed field point $(r,\theta)$
and a fixed particle position, all the quantities of Sec.~\ref{sec:field}
depend on the azimuth only through $\psi$, and only through the three
functions $\sigma$, $\cos(\psi/2)$ and $\sin\psi$. The separations $\delta r$
and $\delta\theta$, the particle data and the operator coefficients are
constants for the azimuthal integral.

Each polynomial that appears in Eqs.~\eqref{eq:first} and \eqref{eq:second},
namely $A$ and its partial derivatives with respect to $\delta r$,
$\delta\theta$, $\sigma$, $\sin\psi$ and $t$, can therefore be collected as
\begin{equation}
  \label{eq:collect}
  \sum_k a_k\,\sigma^k + \sin\psi\sum_k b_k\,\sigma^k ,
\end{equation}
where the coefficients $a_k$ and $b_k$ are polynomials in $\delta r$ and
$\delta\theta$. Every such polynomial is at most linear in $\sin\psi$, because
$A$ is. For later use we write the two parts of Eq.~\eqref{eq:collect} as
$X^{(0)}(u)$ and $X^{(1)}(u)$, so that a polynomial reads
$X^{(0)}(u)+\sin\psi\,X^{(1)}(u)$. Where the powers of $\sigma$ are all odd,
as for $\partial_\sigma A$, we take the common factor $\sigma$ out of both
parts.

\subsection{Clearing the denominators of the source}
\label{sec:clear}

Equations~\eqref{eq:first} and \eqref{eq:second} have denominators $s^9$ and
$s^{11}$, and the source contains derivatives up to second order. Multiplying
the source by $s^{11}$ therefore removes every denominator,
\begin{equation}
  \label{eq:numerator}
  \begin{split}
  N \equiv s^{11}S ={}& \sum_X b_X\,s^2\Bigl(s^2\,\partial_X A - \tfrac72A\,\partial_X s^2\Bigr) \\
    &+ \sum_{X,Y}\frac{b_{XY}}{4}\Bigl[63A\,\partial_X s^2\,\partial_Y s^2
    - 14s^2\bigl(\cdots\bigr) + 4s^4\,\partial_X\partial_Y A\Bigr],
  \end{split}
\end{equation}
where the bracket $(\cdots)$ is that of Eq.~\eqref{eq:second}. With
$\sin\psi=2\sigma\cos(\psi/2)$ and $\cos\psi=1-2u$, the numerator $N$ is a
polynomial in $\sigma$, $\cos(\psi/2)$ and $s^2$, and it is linear in the
coefficients $a_k$ and $b_k$ of the polynomials of Eq.~\eqref{eq:collect}.

Two identities reduce $N$ further. First, $s^2=\alpha+\beta u$ is itself a
polynomial in $u$. Second, $\cos^2(\psi/2)=1-u$, so every even power of
$\cos(\psi/2)$ is also a polynomial in $u$. After these substitutions each
monomial of $N$ is $\sigma^a\cos^b(\psi/2)$ with $b$ equal to zero or one.

The parity of the remaining monomials is fixed by periodicity. Both $\sigma$
and $\cos(\psi/2)$ change sign when $\psi$ increases by $2\pi$, while $N$ is
$2\pi$-periodic, so $a+b$ is even in every monomial. If $b=0$, then $a$ is
even and the monomial is a power of $u$. If $b=1$, then $a$ is odd and the
monomial is $\sigma\cos(\psi/2)=\tfrac12\sin\psi$ times a power of $u$. The
numerator therefore takes the form
\begin{equation}
  \label{eq:PQ}
  N = P(u) + \tfrac12\sin\psi\;Q(u), \qquad
  P(u) = \sum_{j=0}^{7} P_j\,u^j, \qquad Q(u) = \sum_{j=0}^{6} Q_j\,u^j ,
\end{equation}
in which $P(u)$ is even in $\psi$ and $\tfrac12\sin\psi\,Q(u)$ is odd. The
effective source is
\begin{equation}
  \label{eq:rational}
  S = \frac{P(u) + \tfrac12\sin\psi\;Q(u)}{s^{11}} .
\end{equation}

Each coefficient $P_j$ and $Q_j$ is a linear combination of the coefficients
$a_k$ and $b_k$, with weights that depend on $\alpha$, $\beta$, $c$, the other
particle data and the operator coefficients $b_X$ and $b_{XY}$. The weights
are independent of $\psi$. All the cancellation between the singular terms of
Eq.~\eqref{eq:source} now takes place inside the 15 numbers $P_j$ and $Q_j$,
which are computed once per field point. The azimuthal dependence is contained
in Eq.~\eqref{eq:rational}, a smooth function of $\psi$ whose azimuthal
integral has a closed form.

\subsection{The singular field and its first derivatives}
\label{sec:fields}

The other quantities have the same structure, with a different power of $s$ in
the denominator. The singular field itself is
\begin{equation}
  \label{eq:PhiS_numerator}
  \PhiS = \frac{A^{(0)}(u) + \tfrac12\sin\psi\,\bigl[2A^{(1)}(u)\bigr]}{s^7}.
\end{equation}
The $\theta$ derivative follows from Eq.~\eqref{eq:first} with
$\partial_\theta s^2 = 2\alpha_{02}\,\delta\theta$,
\begin{equation}
  \label{eq:dth_numerator}
  \partial_\theta\PhiS = \frac{s^2\,\partial_{\delta\theta}A - 7\alpha_{02}\,\delta\theta\,A}{s^9},
\end{equation}
where $\partial_{\delta\theta}A$ is collected in the form of
Eq.~\eqref{eq:collect}.

The $r$ derivative needs more care, because $\psi$ depends on $r$ through
Eq.~\eqref{eq:psi}. The total derivatives of $A$ and $s^2$ with respect to
$r$ are
\begin{equation}
  \label{eq:total_r}
  D_rA = \partial_{\delta r}A + \partial_\sigma A\,\partial_r\sigma + \partial_{\sin\psi}A\,\partial_r\sin\psi,
  \qquad
  D_r s^2 = 2\alpha_{20}\,\delta r - \tfrac12 c\beta\sin\psi,
\end{equation}
with $\partial_r\sigma=-\tfrac{c}{2}\cos(\psi/2)$ and
$\partial_r\sin\psi=-c\cos\psi=-c(1-2u)$. Inserting the collected forms and
using $\sigma^2\cos^2(\psi/2)=u(1-u)$ gives
$D_rA = D_{\rm e}(u) + \tfrac12\sin\psi\,D_{\rm o}(u)$ with
\begin{align}
  D_{\rm e} &= (\partial_{\delta r}A)^{(0)} - c\,u(1-u)\,(\partial_\sigma A)^{(1)} - c(1-2u)\,(\partial_{\sin\psi}A)^{(0)}, \label{eq:De}\\
  D_{\rm o} &= 2(\partial_{\delta r}A)^{(1)} - \tfrac12 c\,(\partial_\sigma A)^{(0)}, \label{eq:Do}
\end{align}
where the parts of $\partial_\sigma A$ are taken without their common factor of
$\sigma$. The $r$ derivative of the singular field is then
\begin{equation}
  \label{eq:dr_numerator}
  \partial_r\PhiS = \frac{s^2\,D_rA - \tfrac72 A\,D_r s^2}{s^9}.
\end{equation}
When the numerator is expanded, the product of the two odd factors
$\sin\psi$ in $A\,D_rs^2$ becomes $\sin^2\psi=4u(1-u)$, so the numerator again
has the form of Eq.~\eqref{eq:PQ}.

\subsection{The common form}

All four quantities therefore have the form
\begin{equation}
  \label{eq:general}
  f = \frac{P(u) + \tfrac12\sin\psi\;Q(u)}{(\alpha+\beta u)^{\nu}}, \qquad
  \nu=\tfrac{11}{2}\ \text{for}\ S,\quad \nu=\tfrac72\ \text{for}\ \PhiS,\quad
  \nu=\tfrac92\ \text{for}\ \partial_r\PhiS\ \text{and}\ \partial_\theta\PhiS,
\end{equation}
with at most eight coefficients in $P$ and seven in $Q$. The remainder of
Part~\ref{part:theory} computes the azimuthal modes of a function of this
form.

\section{Azimuthal integration}
\label{sec:integration}

\subsection{Phase and parity}

At a fixed field point, Eq.~\eqref{eq:psi} gives $\phi=\phi_0+\psi$ with the
constant $\phi_0=\phi_p+c\,\delta r$. Because $f$ is $2\pi$-periodic in $\psi$,
we can integrate over $\psi$ instead of $\phi$ and write
$e^{-\ii m\phi}=e^{-\ii m\phi_0}\bigl[\cos(m\psi)-\ii\sin(m\psi)\bigr]$. The
even part of $f$ contributes only with $\cos(m\psi)$ and the odd part only with
$\sin(m\psi)$, so
\begin{equation}
  \label{eq:phase}
  f_m = e^{-\ii m\phi_0}\bigl(E_m - \ii\,O_m\bigr),
\end{equation}
with the real integrals
\begin{equation}
  \label{eq:even_odd}
  E_m = \int_0^{2\pi}\frac{\cos(m\psi)\,P(u)}{s^{2\nu}}\,\dd\psi, \qquad
  O_m = \frac12\int_0^{2\pi}\frac{\sin(m\psi)\,\sin\psi\;Q(u)}{s^{2\nu}}\,\dd\psi .
\end{equation}
Since $f$ is real, the mode $-m$ is the complex conjugate of the mode $m$, and
we only need $m\ge0$.

\subsection{The kernels}

Both integrals can be reduced to the family of kernels
\begin{equation}
  \label{eq:kernel}
  \Kmn{\mu}{\nu}(\alpha,\beta) = \int_0^{2\pi}\frac{\cos(\mu\psi)}{(\alpha+\beta\sin^2\frac{\psi}{2})^{\nu}}\,\dd\psi
  = \int_0^{2\pi}\frac{\cos(\mu\psi)}{(A_0 - B_0\cos\psi)^{\nu}}\,\dd\psi,
\end{equation}
for integer $\mu\ge0$ and half-integer $\nu$, where the second form uses
$\sin^2(\psi/2)=(1-\cos\psi)/2$ with $A_0=\alpha+\beta/2$ and $B_0=\beta/2$.
The kernels are the azimuthal modes of a single power of the distance function.
For positive $\nu$ they are positive and decrease with $\mu$, as the binomial
expansion of $(1-x\cos\psi)^{-\nu}$ in powers of $\cos\psi$ shows. Their
evaluation is the subject of Sec.~\ref{sec:kernels}.

To reduce Eq.~\eqref{eq:even_odd} to kernels, the polynomials $P$ and $Q$ must
be expressed in a basis whose products with $\cos(m\psi)$ or
$\sin(m\psi)\sin\psi$ integrate to kernels. There are two natural choices.
Both are exact, and they differ in which kernels they need and in how
rounding errors propagate.

\subsection{The \texorpdfstring{$s^2$}{s squared} basis}

The substitution $u=(s^2-\alpha)/\beta$ writes $P$ as a polynomial in $s^2$,
\begin{equation}
  \label{eq:s2basis}
  P(u) = \sum_k e_k\,s^{2k}, \qquad\text{so that}\qquad E_m = \sum_k e_k\,\Kmn{m}{\nu-k}.
\end{equation}
Each power $s^{2k}$ lowers the effective exponent of the denominator from $\nu$
to $\nu-k$. The coefficients $e_k$ follow from the $P_j$ by scaling the
coefficient of $u^j$ by $\beta^{-j}$ and then shifting the argument of the
polynomial by $-\alpha$.

The odd integral needs one more step, because $\sin(m\psi)\sin\psi$ is not a
cosine. We use the identity
\begin{equation}
  \label{eq:chebyshev}
  \frac{\sin(m\psi)}{\sin\psi}=\sum_{l=0}^{m-1}\cos\bigl((m-1-2l)\psi\bigr)
\end{equation}
and write the remaining factor as
$\tfrac12\sin^2\psi\,Q(u) = 2u(1-u)Q(u)\equiv G(u)$, a polynomial of two
degrees higher than $Q$. With $G(u)=\sum_k g_k\,s^{2k}$,
\begin{equation}
  \label{eq:odd_s2}
  O_m = \sum_k g_k\,L_m^{(\nu-k)}, \qquad
  L_m^{(\nu)} = \sum_{l=0}^{m-1}\Kmn{|m-1-2l|}{\nu} .
\end{equation}
The partial sums $L_m$ need not be formed term by term. They satisfy
$L_0=0$, $L_1=\Kmn{0}{\nu}$ and $L_m=L_{m-2}+2\Kmn{m-1}{\nu}$.

For the source, $P$ has degree seven and $G$ degree eight, so the $s^2$ basis
needs the kernels of the nine orders $\nu=11/2,9/2,\dots,-5/2$, for modes
$\mu\le m$. The kernels of negative order are integrals of positive powers of
the distance function and remain finite as $C_1\to0$. We find numerically that
the condition number of the contraction stays below about $3\times10^3$ for
$C_1<0.03$ and $m\le40$. For larger separations it grows, and we use the
second basis instead.

\subsection{The Fourier basis}

Alternatively we can expand the numerator in a Fourier series of $\psi$,
$P(u)=\sum_l p_l\cos(l\psi)$ and $\tfrac12\sin\psi\,Q(u)=\sum_l q_l\sin(l\psi)$.
The series are finite, since $u^j=\sin^{2j}(\psi/2)=\bigl[(1-\cos\psi)/2\bigr]^j$
has the expansion
\begin{equation}
  \label{eq:fourier_u}
  u^j = 4^{-j}\binom{2j}{j} + 2\sum_{l=1}^{j}(-1)^l\,4^{-j}\binom{2j}{j-l}\cos(l\psi),
\end{equation}
and $\sin\psi\cos(l\psi)$ is a combination of $\sin((l\pm1)\psi)$. The
products $\cos(m\psi)\cos(l\psi)$ and $\sin(m\psi)\sin(l\psi)$ are sums of
cosines of $(m\pm l)\psi$, which gives
\begin{equation}
  \label{eq:fourier_contraction}
  E_m = \sum_{l\ge0} p_l\,\frac{\Kmn{m+l}{\nu}+\Kmn{|m-l|}{\nu}}{2}, \qquad
  O_m = \sum_{l\ge1} q_l\,\frac{\Kmn{|m-l|}{\nu}-\Kmn{m+l}{\nu}}{2}.
\end{equation}
This basis needs a single order $\nu$ but modes up to $\mu=m+7$.

The coefficients of Eq.~\eqref{eq:fourier_u} alternate in sign. Near the
particle the kernels vary slowly with $\mu$, on a scale of order
$1/\sqrt{C_1}$, so the sums in Eq.~\eqref{eq:fourier_contraction} combine
nearly equal kernels with alternating signs and cancel. We find that the
condition number of this contraction grows like $10^8\,(0.044/C_1)^5$, which is
about $10^8$ at $C_1=0.04$ and reaches about $10^{14}$ at
$C_1=3\times10^{-3}$. For larger separations it is small.

\subsection{Choosing the basis}
\label{sec:choice}

The two bases are complementary. The $s^2$ basis is well conditioned near the
particle and the Fourier basis away from it, and we switch between them at
$C_1=0.03$. A second consideration is the cost and accuracy of the kernels
themselves (Sec.~\ref{sec:mu}). The kernels of the $s^2$ basis are cheapest
when a forward recurrence in $\mu$ is accurate up to $\mu=m$. Between
$C_1=3\times10^{-3}$ and $0.03$ this fails for $m$ above about 20, and there we
also use the Fourier basis. Its condition number in that range is at most about
$10^{14}$, which the double-double arithmetic of Sec.~\ref{sec:precision}
absorbs. Below $C_1=3\times10^{-3}$ the Fourier condition number is larger
still, and we keep the $s^2$ basis and compute its kernels by the slower
Miller recurrence when needed. Since all four quantities of
Eq.~\eqref{eq:general} share $\alpha$ and $\beta$, a single table of kernels
serves all of them.

\section{Evaluation of the kernels}
\label{sec:kernels}

\subsection{The zero mode}

For $\mu=0$ and $\nu=\pm1/2$ the kernels are complete elliptic integrals. With
$t=\psi/2$, the substitution $t\to\pi/2-t$ and
$\alpha+\beta\cos^2t=(\alpha+\beta)(1-k^2\sin^2t)$ one finds
\begin{equation}
  \label{eq:elliptic}
  \Kmn{0}{1/2} = \frac{4K(k)}{\sqrt{\alpha+\beta}}, \qquad
  \Kmn{0}{-1/2} = 4\sqrt{\alpha+\beta}\;E(k), \qquad k^2=\frac{\beta}{\alpha+\beta},
\end{equation}
where $K$ and $E$ are the complete elliptic integrals of the first and second
kind. Near the particle $k\to1$, and $K(k)$ diverges like
$\tfrac12\ln\bigl(16/(1-k^2)\bigr)$. Computing $K$ and $E$ by the
arithmetic--geometric mean from the complementary parameter
$1-k^2=\alpha/(\alpha+\beta)$, which is small but known to full relative
precision, avoids the loss of digits that forming $k^2$ first would cause.

The other orders of the zero mode follow from the three-term relation
\begin{equation}
  \label{eq:nu_recurrence}
  p\,D\,\Kmn{0}{p+1} = (2p-1)\,A_0\,\Kmn{0}{p} - (p-1)\,\Kmn{0}{p-1}, \qquad
  D = A_0^2-B_0^2 = \alpha(\alpha+\beta),
\end{equation}
applied upward from $\nu=1/2$ to $11/2$ and downward from $\nu=-1/2$ to
$-7/2$. In both directions the kernels are the dominant solution of the
recurrence, so rounding errors do not grow. The first mode of each order
follows from the zero modes of two neighbouring orders,
\begin{equation}
  \label{eq:K1}
  \Kmn{1}{\nu} = \frac{A_0\,\Kmn{0}{\nu}-\Kmn{0}{\nu-1}}{B_0},
\end{equation}
which is Eq.~\eqref{eq:kernel} with $\cos\psi=(A_0-s^2)/B_0$.

\subsection{The recurrence in \texorpdfstring{$\mu$}{mu}}
\label{sec:mu}

The kernels of a fixed order satisfy the three-term recurrence
\begin{equation}
  \label{eq:mu_recurrence}
  (\mu+1-\nu)\,\Kmn{\mu+1}{\nu} - 2\mu z\,\Kmn{\mu}{\nu} + (\mu+\nu-1)\,\Kmn{\mu-1}{\nu} = 0,
  \qquad z=\frac{A_0}{B_0}=1+2C_1 .
\end{equation}
For large $\mu$ its two independent solutions behave like $\rho^{\pm\mu}$, up to
powers of $\mu$, with
\begin{equation}
  \label{eq:rho}
  \rho=z+\sqrt{z^2-1}, \qquad \ln\rho = \mathrm{arcosh}(1+2C_1) = 2\,\mathrm{arsinh}\sqrt{C_1}.
\end{equation}
The kernels are the minimal solution, decreasing like $\rho^{-\mu}$. The rate
$\ln\rho$ equals the distance of the singularities of the integrand from the
real axis found in Sec.~\ref{sec:direct}, as expected for the Fourier
coefficients of an analytic periodic function.

Running Eq.~\eqref{eq:mu_recurrence} forward from $\Kmn{0}{\nu}$ and
$\Kmn{1}{\nu}$ is the cheapest way to obtain the kernels, but any rounding
error excites the dominant solution, which grows relative to the kernels like
$\rho^{2\mu}$. Near the particle $\rho\to1$, and the forward recurrence is
accurate to high $\mu$. Further away, or for large $m$, we predict the error
at the largest required $\mu$ from a fit to a recurrence run at 32 digits, and
accept the forward direction when the prediction is small enough.

Otherwise we use Miller's algorithm. Starting from arbitrary values, for
example zero and one, at a guard order well above the largest required $\mu$,
we run the recurrence downward. In this direction the minimal solution grows
and the dominant solution decays, so after a number of steps set by
$\ln\rho$ the iterates are proportional to the kernels to working precision.
The overall factor follows from the known value of $\Kmn{0}{\nu}$.

The relations~\eqref{eq:elliptic}--\eqref{eq:mu_recurrence} agree with 30-digit
quadrature of Eq.~\eqref{eq:kernel} to better than $10^{-24}$ relative for the
orders and modes we tested.

\section{Precision}
\label{sec:precision}

\subsection{Rounding in the integration}

The changes of basis, the kernels and the contractions of
Sec.~\ref{sec:integration} are exact algebra applied to the coefficients of
Eq.~\eqref{eq:general}. Their rounding errors are multiplied by the condition
numbers discussed above, which reach $3\times10^3$ in the $s^2$ basis and
$10^{14}$ in the Fourier basis at the edge of its range. We therefore carry out
these steps in double-double arithmetic, in which a number is the unevaluated
sum of two doubles and holds about 106 bits, and round the result to double
once. The modes are then the integrals of the given coefficients to about
$10^{-16}$ relative to each mode, and the choice of basis and recurrence only
affects the run time.

\subsection{Cancellation in the numerator coefficients}

The precision of the coefficients $P_j$ and $Q_j$ is a separate question. Near
the particle they are the small remainder of products as large as the singular
terms of Eq.~\eqref{eq:source}. A rounding error in one of these coefficients
perturbs the numerator by the same polynomial at every azimuth. Unlike the
rounding of a quadrature node, it does not average out in the integral. A
relative rounding error $\epsilon$ in the products therefore changes $S_m$ by
about $\epsilon$ times the azimuthal integral of the singular terms, which grows
like $1/C_1$ (Sec.~\ref{sec:termwise}), while $S_m$ stays finite. When the
coefficients are computed in double precision, the source modes have errors
between about $2\times10^{-15}/C_1$ and $10^{-13}/C_1$ relative to the largest
mode, which we measured from $C_1=10^{-1}$ to $10^{-10}$.

We therefore compute the coefficients in double-double arithmetic below a
threshold in $C_1$ and in double precision above it. Double-double coefficients
reproduce the exact modes of the given particle data to about $10^{-16}$. The
threshold balances run time against accuracy, and we discuss its value in
Sec.~\ref{sec:prepare}.

\subsection{Rounding of the particle data}

The particle data are themselves computed and stored in double precision.
Perturbing each of them by one unit in the last place changes the source modes
by about $5.6\times10^{-15}/C_1$ relative to the largest mode. This
contribution is set by the precision of the particle data, and no choice of
arithmetic in the steps above can reduce it. It is the leading error of the
method near the particle.

\part{Implementation}
\label{part:implementation}

\section{Overview}
\label{sec:overview}

The method is implemented in three layers, which correspond to the stages of
Part~\ref{part:theory} and run at different frequencies.
\begin{enumerate}
  \item \code{effsource\_equatorial\_ctx\_set\_particle}, in
    \code{kerr\-equatorial\-coeffs\-ctx.c}, computes the particle data of
    Sec.~\ref{sec:field} once per particle position and stores them in
    \code{struct effsource\_equatorial\_ctx}. It calls
    \code{effsource\_equatorial\_ctx\_set\_particle\_dt} and
    \code{\_dtt}, in \code{kerr\-equatorial\-dtcoeffs\-ctx.c} and
    \code{kerr\-equatorial\-dttcoeffs\-ctx.c}, for the time derivatives.
  \item \code{effsource\_equatorial\_ctx\_prepare} computes the coefficients of
    Sec.~\ref{sec:rational} once per field point and particle position, and
    stores them in \code{struct effsource\_equatorial\_point}. It is generated
    code, in \code{kerr\-equatorial\-point\-ctx.c}.
  \item \code{effsource\_equatorial\_ctx\_analytic\_modes}, in
    \code{kerr\-equatorial\-mmode.c}, builds the numerators of
    Eq.~\eqref{eq:general}, computes the kernels and returns the modes
    $m=0,\dots,m_{\max}$ of any of the four quantities.
\end{enumerate}
The pointwise singular field and source are computed by
\code{effsource\_equatorial\_ctx\_calc} in \code{kerr\-equatorial\-ctx.c}, which
is the source from which the code generator works. The term-by-term route of
Sec.~\ref{sec:termwise} is \code{effsource\_equatorial\_ctx\_calc\_m} in the
same file, with tabulated closed forms for $m\le20$, and Sec.~\ref{sec:old}
compares it with the method of this note.

In SpECTRE, \code{EccentricOrbit\_detail::m\_modes} calls the second and third
layers at every grid point and time sample with $m_{\max}=|m|$, and takes the
complex conjugate for negative $m$.

\section{Code generation}
\label{sec:codegen}

The script \code{codegen/point\_codegen.py} reads \code{kerr-equatorial-ctx.c}
and writes \code{effsource\_equatorial\_point.h} and
\code{kerr\-equatorial\-point\-ctx.c}. It is run from the repository root as
\begin{verbatim}
python codegen/point_codegen.py kerr-equatorial-ctx.c \
    effsource_equatorial_point.h kerr-equatorial-point-ctx.c
\end{verbatim}
and performs the steps of Sec.~\ref{sec:rational} symbolically with SymPy.
\begin{enumerate}
  \item It parses the body of \code{effsource\_equatorial\_ctx\_calc}, which is
    straight-line C, and splits it at the comment that introduces the total
    derivatives. The first part defines $A$ and its partial derivatives.
  \item It collects each of these 16 polynomials in powers of $\sigma$
    (\code{dQ} in the code) and $\sin\psi$ (\code{dR}), checks that each is
    linear in $\sin\psi$, and lays out each of the resulting 28 parts as a
    Horner polynomial in $u$, with an overall factor $\sigma$ where the powers
    are odd. These parts are the fields of \code{struct
    effsource\_equatorial\_point}, with 121 coefficients in total. The suffixes
    \code{\_R0} and \code{\_R1} of the field names denote $X^{(0)}$ and
    $X^{(1)}$.
  \item It executes the whole calc body symbolically, with the trigonometric
    functions of $\psi$ replaced by $\sigma$ and $\cos(\psi/2)$ and with
    $\sqrt{s^2}$ replaced by $(s^2)^{-5}$. This substitution maps each power
    $s^{-(2k+1)}$ of the calc body to $(s^2)^{5-k}=s^{11-(2k+1)}$, so the
    symbolic source is the numerator $N$ of Eq.~\eqref{eq:numerator}, without
    division.
  \item It expands $N$ as a polynomial in $\sigma$, $\cos(\psi/2)$ and $s^2$,
    keeping the scalar factors unexpanded, substitutes $s^2=\alpha+\beta u$ and
    $\cos^2(\psi/2)=1-u$, and checks that every remaining monomial has the
    parity of Eq.~\eqref{eq:PQ}. The result expresses $P_0,\dots,P_7$ and
    $Q_0,\dots,Q_6$ as sums of an integer times one of the 121 coefficients
    times one of 65 weights, 547 products in total.
  \item It writes the coefficient polynomials, the weights after common
    subexpression elimination, and the products twice: once in double-double
    arithmetic as \code{effsource\_equatorial\_ctx\_prepare\_dd} and once in
    double as \code{effsource\_equatorial\_ctx\_prepare\_double}.
  \item It also writes two pointwise evaluators for tests:
    \code{effsource\_equatorial\_ctx\_sweep\_calc}, the calc body with every
    polynomial replaced by its Horner form, and
    \code{effsource\_equatorial\_ctx\_sweep\_src}, which evaluates
    Eq.~\eqref{eq:rational} at a given azimuth.
\end{enumerate}
The generated files should not be edited by hand. A change to the calc body or
to the threshold of Sec.~\ref{sec:prepare} is made in the source or in the
script, followed by a regeneration.

\section{The prepare functions}
\label{sec:prepare}

The double-double arithmetic is in \code{effsource\_dd.h}. It follows Hida, Li
and Bailey and uses only IEEE double operations, with fused multiply--add where
the compiler reports it fast and Dekker's splitting otherwise, so the results
are the same on every platform.

\code{effsource\_equatorial\_ctx\_prepare\_dd} forms $\delta r$ and
$\delta\theta$ by exact two-sums, evaluates $\sin\theta$, $\sin2\theta$ and
$\cos2\theta$ from the double-double sine and cosine of $\theta-\pi/2$, and
evaluates the 121 coefficients, the 65 weights and the 547 products in
double-double. It stores the high parts of all coefficients. For $P_j$, $Q_j$,
$\alpha$ and the nine polynomial parts that the numerators of
Sec.~\ref{sec:fields} use, it also stores the low parts in the fields ending in
\code{\_lo}, which the m-mode routine reads.
\code{effsource\_equatorial\_ctx\_prepare\_double} evaluates the same
expressions in double and sets the low parts to zero. With double coefficients,
the coefficient rounding is amplified by the condition number of the
contraction. This limits the accuracy of small high-$m$ modes relative to
themselves, to about $10^{-10}$ at $C_1\approx0.05$ and $m=20$, while the
error relative to the largest mode remains near $10^{-15}$.

\code{effsource\_equatorial\_ctx\_prepare} chooses between the two by the
threshold \code{EFFSOURCE\_PREPARE\_DD\_C1} $=10^{-3}$, using double-double for
$C_1<10^{-3}$. With clang at \code{-O2} the double-double version takes about
$11\,\mu$s per field point and the double version $0.64\,\mu$s. Above the
threshold the source modes are accurate to about $10^{-10}$ relative to the
largest mode, with a largest measured error of $6\times10^{-11}$ at
$C_1=10^{-3}$. On a representative SpECTRE grid near an orbit with $e=0.2$,
about $20\%$ of the (grid point, time) pairs fall below the threshold,
compared with $57\%$ for the earlier threshold of $10^{-2}$. In tests over 13
eccentric orbits with $m\le40$ and $|n|\le50$, the radial modes computed with
the threshold at $10^{-3}$ differ from those with double-double everywhere by
at most $4\times10^{-10}$ times the mean modulus of the $m$-modes in their
trapezoidal sum over time.

\section{The m-mode routine}
\label{sec:mmode}

\subsection{Interface}

The routine
\begin{verbatim}
int effsource_equatorial_ctx_analytic_modes(ctx, point, m_max,
                                            src, PhiS, dPhiS_dr, dPhiS_dth);
\end{verbatim}
writes the modes $m=0,\dots,m_{\max}$ of each quantity whose output pointer is
not \code{NULL}, with the real and imaginary parts of mode $m$ at indices $2m$
and $2m+1$, the layout of an array of complex doubles. All outputs share one
table of kernels, so each output after the first adds only its contraction.
The routine returns 1 if $m_{\max}+7$ exceeds
\code{EFFSOURCE\_MMODE\_MAX\_MU} $=256$, which limits the modes to
$m\le249$, and 2 if $\alpha\le0$, which happens only on the orbit. The
convention is that of Eq.~\eqref{eq:mode}, $f_m=\int f\,e^{-\ii m\phi}\,\dd\phi$.
The wrapper \code{effsource\_equatorial\_ctx\_analytic\_src\_modes} returns the
source modes as separate real and imaginary arrays.

\subsection{Numerators}

The numerators are assembled in double-double from the high and low parts of
the point structure. \code{src\_numerator} loads $P_j$ and $Q_j$ from it.
\code{PhiS\_numerator}, \code{dPhiS\_dth\_numerator} and
\code{dPhiS\_dr\_numerator} implement
Eqs.~\eqref{eq:PhiS_numerator}--\eqref{eq:dr_numerator} from the fields
\code{A\_R0}, \code{A\_R1}, \code{dA\_dr\_R0}, \code{dA\_dr\_R1},
\code{dA\_dth\_R0}, \code{dA\_dth\_R1}, \code{dA\_dQ\_R0}, \code{dA\_dQ\_R1} and
\code{dA\_dR\_R0}.

\subsection{Kernels}

\code{zero\_mode\_kernels} computes $\Kmn{0}{\nu}$ for
$\nu=-7/2,\dots,11/2$ from Eqs.~\eqref{eq:elliptic} and
\eqref{eq:nu_recurrence}, with the arithmetic--geometric mean of
\code{elliptic\_KE}. \code{kernels\_dd} then fills the table
$\Kmn{\mu}{\nu}$ for the orders and modes the chosen basis needs. The function
\code{forward\_is\_accurate} predicts the error of the forward recurrence at
$\mu_{\max}$ as
\begin{equation}
  \label{eq:forward_fit}
  \log_{10}\varepsilon \approx -30 + \frac{2\mu_{\max}\ln\rho}{\ln10}
    + \max\bigl(0,\;7.5\log_{10}\mu_{\max}-5.5\bigr),
\end{equation}
a fit to a 32-digit recurrence for $\nu=11/2,\dots,-5/2$ and $C_1$ from
$10^{-3}$ to $0.3$, and accepts the forward direction when this is below
$10^{-19}$ (\code{FORWARD\_LOG10\_ERROR}). Otherwise the Miller recurrence
starts 20 orders above the order at which the dominant solution has decayed by
$10^{-34}$ (\code{MILLER\_LOG10\_START}), and rescales by $2^{-600}$ whenever an
iterate exceeds $2^{600}$ to avoid overflow. The public function
\code{effsource\_equatorial\_mmode\_kernels} returns the kernel table rounded
to double.

\subsection{Basis choice and contractions}

The function \code{modes} selects the $s^2$ basis when $C_1<0.03$
(\code{EFFSOURCE\_MMODE\_FOURIER\_C1}) and either $C_1<3\times10^{-3}$
(\code{FOURIER\_MIN\_C1}) or the forward recurrence is accurate up to
$\mu=m_{\max}$, as described in Sec.~\ref{sec:choice}. In the
$s^2$ basis it needs the orders from $11/2$ down to $\nu-\deg G$ for each output
and modes up to $m_{\max}$. In the Fourier basis it needs the order of each
output and modes up to $m_{\max}+7$. \code{s2\_contraction} performs the
changes of basis with \code{to\_s2\_basis}, a scaling followed by a Taylor
shift, for $P$ and $G$, and evaluates Eqs.~\eqref{eq:s2basis} and
\eqref{eq:odd_s2} with the partial sums $L_m$ updated by their recurrence.
\code{fourier\_contraction} forms the Fourier coefficients from
Eq.~\eqref{eq:fourier_u} and applies Eq.~\eqref{eq:fourier_contraction}.
Finally \code{modes} applies the phase of Eq.~\eqref{eq:phase}, with the angle
$m\phi_0$ formed as an exact product in double-double and the cosine and sine
corrected for its low part.

\section{The earlier implementation}
\label{sec:old}

The term-by-term route of Sec.~\ref{sec:termwise} is implemented in
\code{effsource\_equatorial\_ctx\_calc\_m}, in \code{kerr\-equatorial\-ctx.c},
and in \code{effsource\_equatorial\_ctx\_PhiS\_m} for the singular field
alone. The coefficients $R_j$ and $I_j$ of Eq.~\eqref{eq:A_split} are the
arrays \code{ReA} and \code{ImA}, which the function evaluates together with
their first and second derivatives in $t$, $r$ and $\theta$. The polynomials of
Eq.~\eqref{eq:old_closed} are tabulated in \code{ReEI} and \code{ImEI}, each of
dimension $21\times2\times5\times27$: the modes $m=0,\dots,20$, the two
elliptic integrals, the five values of $j$ and the powers of $C_1$ up to 26.
For $m>20$ the function prints a message and returns without output. By
default the elliptic integrals come from GSL as functions of $k$. Setting the
flag \code{ell\_from\_complement} in the particle context selects the
arithmetic--geometric mean from the complementary parameter instead.

Table~\ref{tab:old} compares \code{effsource\_equatorial\_ctx\_calc\_m} with
the analytic modes from double-double coefficients, which are accurate to
about $10^{-14}$ relative to each mode for the given particle data
(Sec.~\ref{sec:validation}). Both use the same particle data, for a
particle at $t=0.1\,T_r$ on an orbit with $a=0.9$, $p=10$ and $e=0.2$, and
field points along one direction in the $(\delta r,\delta\theta)$ plane. The
differences are therefore the errors of the term-by-term route. The table uses
the complementary parameter for the elliptic integrals.

\begin{table}[H]
  \centering
  \caption{Relative error of each mode of \code{effsource\_equatorial\_ctx\_calc\_m},
    with \code{ell\_from\_complement} set, against the analytic modes. The last
    column is the largest source error over $m\le20$ relative to the largest
    source mode.}
  \label{tab:old}
  \begin{tabular}{rrrrrrrr}
    \toprule
    & \multicolumn{3}{c}{$\PhiS$} & \multicolumn{3}{c}{$S$} & $S$, largest \\
    \cmidrule(lr){2-4}\cmidrule(lr){5-7}
    $C_1$ & $m=0$ & $m=10$ & $m=20$ & $m=0$ & $m=10$ & $m=20$ & mode \\
    \midrule
    $1$       & $3\times10^{-16}$ & $7$                 & $10^{15}$           & $2\times10^{-14}$ & $2$                 & $10^{15}$          & $2\times10^{5}$ \\
    $10^{-1}$ & $3\times10^{-16}$ & $4\times10^{-10}$   & $8\times10^{-4}$    & $3\times10^{-13}$ & $4\times10^{-7}$    & $3\times10^{-2}$   & $10^{-5}$ \\
    $10^{-2}$ & $3\times10^{-16}$ & $10^{-13}$          & $10^{-11}$          & $3\times10^{-13}$ & $9\times10^{-10}$   & $2\times10^{-7}$   & $10^{-9}$ \\
    $10^{-4}$ & $4\times10^{-16}$ & $2\times10^{-15}$   & $6\times10^{-15}$   & $4\times10^{-9}$  & $8\times10^{-7}$    & $10^{-6}$          & $4\times10^{-9}$ \\
    $10^{-6}$ & $3\times10^{-16}$ & $10^{-16}$          & $2\times10^{-16}$   & $2\times10^{-7}$  & $2\times10^{-5}$    & $8\times10^{-4}$   & $4\times10^{-7}$ \\
    $10^{-8}$ & $2\times10^{-16}$ & $5\times10^{-16}$   & $10^{-15}$          & $8\times10^{-6}$  & $10^{-2}$           & $3\times10^{-2}$   & $10^{-4}$ \\
    \bottomrule
  \end{tabular}
\end{table}

The three failures of Sec.~\ref{sec:termwise} appear in separate regions of
the table. At $C_1=1$ and $10^{-1}$ the closed forms of the high modes cancel,
and the singular field is already wrong at $m=10$ for $C_1=1$, with the source
errors following the singular field. For small $C_1$ the modes of $\PhiS$ are
accurate to $10^{-15}$, but the source errors grow like $1/C_1$, and the
largest source error reaches $10^{-4}$ of the largest mode at $C_1=10^{-8}$. At
intermediate separations, $C_1=10^{-2}$, the high source modes are accurate to
only about $10^{-7}$. With the default elliptic integrals, the modes of $\PhiS$
also have relative errors of $8\times10^{-14}$ at $C_1=10^{-4}$,
$10^{-11}$ at $10^{-6}$ and $2\times10^{-10}$ at $10^{-8}$, while the source
errors are of the same size as in the table.

For comparison, the analytic modes with the default threshold of
Sec.~\ref{sec:prepare} have source errors below about $6\times10^{-11}$ of the
largest mode for $C_1\ge10^{-3}$ and of order $10^{-16}$ below, for the same
particle data.

Table~\ref{tab:old_fields} adds the first derivatives of the singular field,
which enter a solver as jump data on the boundary of the region where the
effective source is used. In the SpECTRE configuration of this orbit, with that
region at $4<r<20$ and $|\cos\theta|<0.5$, $C_1$ on the boundary ranges from
$0.058$, on the polar edges at the radii of the orbit, to $0.59$, on the outer
radial edge with the particle at periapsis.

\begin{table}[H]
  \centering
  \caption{Relative error of each mode of \code{effsource\_equatorial\_ctx\_calc\_m},
    with \code{ell\_from\_complement} set, against the analytic modes, for the
    singular field, its first derivatives and the source.}
  \label{tab:old_fields}
  \small
  \setlength{\tabcolsep}{4pt}
  \begin{tabular}{rrrrrrrrr}
    \toprule
    & \multicolumn{2}{c}{$\PhiS$} & \multicolumn{2}{c}{$\partial_r\PhiS$}
    & \multicolumn{2}{c}{$\partial_\theta\PhiS$} & \multicolumn{2}{c}{$S$} \\
    \cmidrule(lr){2-3}\cmidrule(lr){4-5}\cmidrule(lr){6-7}\cmidrule(lr){8-9}
    $C_1$ & $m=2$ & $m=10$ & $m=2$ & $m=10$ & $m=2$ & $m=10$ & $m=0$ & $m=10$ \\
    \midrule
    $0.7$     & $1\times10^{-13}$ & $0.1$               & $1\times10^{-12}$ & $5$                 & $7\times10^{-13}$ & $0.06$              & $8\times10^{-14}$ & $0.07$ \\
    $0.3$     & $1\times10^{-14}$ & $6\times10^{-6}$    & $9\times10^{-14}$ & $2\times10^{-5}$    & $1\times10^{-13}$ & $9\times10^{-6}$    & $6\times10^{-13}$ & $1\times10^{-4}$ \\
    $0.05$    & $2\times10^{-15}$ & $8\times10^{-11}$   & $6\times10^{-15}$ & $8\times10^{-11}$   & $7\times10^{-15}$ & $1\times10^{-10}$   & $9\times10^{-13}$ & $4\times10^{-9}$ \\
    $10^{-2}$ & $1\times10^{-15}$ & $1\times10^{-13}$   & $1\times10^{-14}$ & $3\times10^{-14}$   & $5\times10^{-15}$ & $8\times10^{-14}$   & $3\times10^{-13}$ & $9\times10^{-10}$ \\
    $10^{-4}$ & $1\times10^{-15}$ & $2\times10^{-15}$   & $3\times10^{-15}$ & $6\times10^{-15}$   & $8\times10^{-15}$ & $1\times10^{-14}$   & $4\times10^{-9}$  & $8\times10^{-7}$ \\
    \bottomrule
  \end{tabular}
\end{table}

The first derivatives behave like the singular field itself. Differentiating
after the projection leaves them accurate to about $10^{-14}$ for small $C_1$,
and they degrade at large $C_1$ and high $m$ together with the closed forms of
$\PhiS$. At $m=10$ and $C_1=0.7$, close to the largest value on the boundary
above, the closed-form $\partial_r\PhiS$ is wrong by a factor of five. At
$m\le2$ the two routes agree to $10^{-12}$ or better at every separation in the
table. Only the source degrades as $C_1$ decreases, already at $m=0$.

\section{Validation and run time}
\label{sec:validation}

The program \code{test/mmode-check.c} makes four comparisons, grouped by
decade of $C_1$ from $10^{-9}$ to $10^{-1}$, and exits with an error if any of
them fails.
\begin{enumerate}
  \item The kernels for $\nu=11/2,\dots,-5/2$ and $\mu\le40$ agree with a
    long-double trapezoidal rule to $1.2\times10^{-15}$ of $\Kmn{0}{\nu}$, for
    $C_1$ from $10^{-8}$ to $10^{3}$.
  \item The source modes for $m\le20$ from the double-double coefficients agree
    with a long-double trapezoidal rule applied to Eq.~\eqref{eq:rational} with
    the same coefficients to $7.6\times10^{-15}$ relative to each mode, up to the
    rounding of the reference. This comparison tests the kernels and
    contractions in isolation.
  \item The source modes from the hybrid coefficients are compared with the
    trapezoidal rule applied to a long-double evaluation of the calc body, for
    $C_1$ from $10^{-5}$ to $10^{-1}$ where the reference grid is small enough
    to use. They agree to $2.4\times10^{-11}$ of the largest mode for
    $10^{-3}\le C_1<10^{-2}$, just above the threshold, and to
    $1.2\times10^{-12}$ or better in the other decades.
  \item The modes of $\PhiS$, $\partial_r\PhiS$ and $\partial_\theta\PhiS$
    agree with quadrature of the long-double chain-rule formulas from the same
    coefficients to $8.9\times10^{-14}$ relative to each mode, and with the
    long-double calc body to $2.3\times10^{-15}$ of the largest mode.
\end{enumerate}
With \code{--dump}, \code{mmode-check} also writes its cases for
\code{test/mmode\_exact\_check.py}, which recomputes each mode with mpmath from
a hypergeometric closed form of the kernels. The program
\code{test/point-check.c} grades \code{sweep\_calc} and \code{sweep\_src}
against the long-double calc body, both pointwise and for the source modes. In
SpECTRE, the unit test of \code{EccentricOrbit} compares the four outputs with
a trapezoidal rule at 12 field points, using \code{sweep\_src} for the source
and the calc body for the singular field, with tolerances that depend on which
prepare function applies.

Table~\ref{tab:timing} lists the run time of the m-mode routine on one core
with clang at \code{-O2}, excluding the prepare function. Including a prepare
call, the source mode $m=2$ takes about $2.3\,\mu$s per field point above the
threshold and about $13\,\mu$s below it.

\begin{table}[H]
  \centering
  \caption{Run time of \code{effsource\_equatorial\_ctx\_analytic\_modes} for
    modes $0$ to $m$ at $C_1\le10^{-2}$, excluding the prepare function.}
  \label{tab:timing}
  \begin{tabular}{rrr}
    \toprule
    $m$ & source only ($\mu$s) & all four outputs ($\mu$s) \\
    \midrule
    0  & 1.3  & 2.9  \\
    2  & 1.7  & 3.8  \\
    10 & 4.2  & 8.3  \\
    20 & 7.2  & 13.9 \\
    40 & 13.3 & 25.2 \\
    \bottomrule
  \end{tabular}
\end{table}

\section{Implementation map}
\label{sec:map}

Table~\ref{tab:map} lists where each step of Part~\ref{part:theory} is
implemented.

\begin{table}[H]
  \centering
  \caption{Where each step is implemented.}
  \label{tab:map}
  \begin{tabular}{>{\raggedright\arraybackslash}p{0.42\textwidth}>{\raggedright\arraybackslash}p{0.5\textwidth}}
    \toprule
    Step & Implementation \\
    \midrule
    Particle data, Sec.~\ref{sec:field}
      & \code{effsource\_equatorial\_ctx\_set\_particle} and its \code{\_dt} and
        \code{\_dtt} companions \\
    Singular field, derivatives and source, Eqs.~\eqref{eq:puncture}--\eqref{eq:source}
      & \code{effsource\_equatorial\_ctx\_calc} in \code{kerr-equatorial-ctx.c} \\
    Collection, Eq.~\eqref{eq:collect}, and the numerator, Eqs.~\eqref{eq:numerator}--\eqref{eq:rational}
      & \code{codegen/point\_codegen.py}, generating \code{kerr-equatorial-point-ctx.c} \\
    Coefficients at a field point and the precision threshold
      & \code{effsource\_equatorial\_ctx\_prepare}, \code{\_prepare\_dd} and
        \code{\_prepare\_double} \\
    Numerators of $\PhiS$ and its derivatives,
      Eqs.~\eqref{eq:PhiS_numerator}--\eqref{eq:dr_numerator}
      & \code{PhiS\_numerator}, \code{dPhiS\_dth\_numerator} and
        \code{dPhiS\_dr\_numerator} in \code{kerr-equatorial-mmode.c} \\
    Kernels, Eqs.~\eqref{eq:elliptic}--\eqref{eq:mu_recurrence}
      & \code{zero\_mode\_kernels}, \code{elliptic\_KE} and \code{kernels\_dd} \\
    Contractions, Eqs.~\eqref{eq:s2basis}--\eqref{eq:fourier_contraction}
      & \code{s2\_contraction} and \code{fourier\_contraction} \\
    Basis choice and phase, Eq.~\eqref{eq:phase}
      & \code{modes} and \code{effsource\_equatorial\_ctx\_analytic\_modes} \\
    Pointwise references for tests
      & \code{effsource\_equatorial\_ctx\_sweep\_src} and \code{\_sweep\_calc} \\
    \bottomrule
  \end{tabular}
\end{table}

\end{document}
